% GENERATED FILE. Edit canonical lessons, then run npm run generate:books.
% canonical-source-hash: fnv1a64:0af416621bad649f
% canonical-lessons: UR-C309-hamla, UR-C309-qila, UR-C309-gherna, UR-C309-top, UR-C309-bhirna

\chapter{The Attack}
\label{ch:ur-the-attack}
\hlchaptermodality{309}

\begin{chapteropening}
\textbf{By the end of this chapter:} \emph{I can tell how an army attacked a fort.}

Complete the last lesson of chapter 309: I can tell how an army attacked a fort.
\end{chapteropening}

\section[\texorpdfstring{\ur{حملہ} (hamlā)}{حملہ (hamlā)}]{\ur{حملہ} (hamlā) — an attack}

\label{lesson:UR-C309-hamla}



\begin{quote}
Before the new one: say the Urdu for a royal court, then the Urdu for an heir.
\end{quote}

\subsection*{You'll want to know: \ur{حملہ}}
\textbf{\ur{حملہ}} — \emph{hamlā} — \textquotedblleft{}an attack\textquotedblright{}.

\begin{quote}
\ur{فوج} \ur{نے} \ur{رات} \ur{کو} \ur{حملہ} \ur{کیا۔}
\end{quote}

\begin{quote}
\emph{fauj ne rāt ko hamlā kiyā.}
\end{quote}

\begin{quote}
The army attacked at night.
\end{quote}

Masculine.

\begin{cousinweb}
\textbf{\ur{حملہ}} (\emph{hamlā}) is usually traced, through Persian, to Arabic \emph{ḥamla}, \textquotedblleft{}a charge, an onset\textquotedblright{}. \textbf{\ur{حملہ} \ur{کیا}} is attacked.
\end{cousinweb}

\begin{grammarlens}[title={What you've built}]
One more word for telling what happened.
\end{grammarlens}

\subsection*{Your turn}
\begin{itemize}
\raggedright
  \item \emph{Say it:} \emph{hamlā}
  \item \emph{Say it:} \emph{hamlā}, once more
  \item \emph{Say it:} \emph{fauj ne rāt ko hamlā kiyā.} — The army attacked at night.
  \item \emph{Recall:} say the Urdu for a royal court, then the Urdu for an heir, then say \emph{hamlā} again
\end{itemize}

\begin{tcolorbox}[breakable,colback=teal!4,colframe=teal!35!black,title={Before you move on}]
What does \ur{حملہ} mean? (\textquotedblleft{}An attack\textquotedblright{}.) Say it once more.
\end{tcolorbox}
\section[\texorpdfstring{\ur{قلعہ} (qilā)}{قلعہ (qilā)}]{\ur{قلعہ} (qilā) — a fort}

\label{lesson:UR-C309-qila}



\begin{quote}
Before the new one: say the Urdu for an heir, then the Urdu for an attack.
\end{quote}

\subsection*{You'll want to know: \ur{قلعہ}}
\textbf{\ur{قلعہ}} — \emph{qilā} — \textquotedblleft{}a fort\textquotedblright{}.

\begin{quote}
\ur{قلعہ} \ur{پہاڑی} \ur{پر} \ur{تھا۔}
\end{quote}

\begin{quote}
\emph{qilā pahāṛī par thā.}
\end{quote}

\begin{quote}
The fort stood on a hill.
\end{quote}

Masculine.

\begin{cousinweb}
\textbf{\ur{قلعہ}} (\emph{qilā}) is usually traced, through Persian, to Arabic \emph{qalʿa}, \textquotedblleft{}a fortress\textquotedblright{}. The oblique is \textbf{\ur{قلعے}}: \textbf{\ur{قلعے} \ur{میں}}.
\end{cousinweb}

\begin{grammarlens}[title={What you've built}]
One more word for telling what happened.
\end{grammarlens}

\subsection*{Your turn}
\begin{itemize}
\raggedright
  \item \emph{Say it:} \emph{qilā}
  \item \emph{Say it:} \emph{qilā}, once more
  \item \emph{Say it:} \emph{qilā pahāṛī par thā.} — The fort stood on a hill.
  \item \emph{Recall:} say the Urdu for an heir, then the Urdu for an attack, then say \emph{qilā} again
\end{itemize}

\begin{tcolorbox}[breakable,colback=teal!4,colframe=teal!35!black,title={Before you move on}]
What does \ur{قلعہ} mean? (\textquotedblleft{}A fort\textquotedblright{}.) Say it once more.
\end{tcolorbox}
\section[\texorpdfstring{\ur{گھیرنا} (ghernā)}{گھیرنا (ghernā)}]{\ur{گھیرنا} (ghernā) — to surround; to lay siege to}

\label{lesson:UR-C309-gherna}



\begin{quote}
Before the new one: say the Urdu for an attack, then the Urdu for a fort.
\end{quote}

\subsection*{You'll want to know: \ur{گھیرنا}}
\textbf{\ur{گھیرنا}} — \emph{ghernā} — \textquotedblleft{}to surround; to lay siege to\textquotedblright{}.

\begin{quote}
\ur{فوج} \ur{نے} \ur{قلعے} \ur{کو} \ur{گھیرا۔}
\end{quote}

\begin{quote}
\emph{fauj ne qile ko gherā.}
\end{quote}

\begin{quote}
The army surrounded the fort.
\end{quote}

Past, with \ur{نے} (he or she …): \ur{اس} \ur{نے} \ur{شہر} \ur{کو} \ur{گھیرا}, \emph{us ne shahr ko gherā}.

\begin{cousinweb}
\textbf{\ur{گھیرنا}} (\emph{ghernā}) is inherited: CDIAL lists the Hindi cousin \emph{ghernā}, \textquotedblleft{}to surround\textquotedblright{}, under a form it reconstructs (CDIAL 4474).
\end{cousinweb}

\begin{grammarlens}[title={What you've built}]
One more word for telling what happened.
\end{grammarlens}

\subsection*{Your turn}
\begin{itemize}
\raggedright
  \item \emph{Say it:} \emph{ghernā}
  \item \emph{Say it:} \emph{ghernā}, once more
  \item \emph{Say it:} \emph{fauj ne qile ko gherā.} — The army surrounded the fort.
  \item \emph{Recall:} say the Urdu for an attack, then the Urdu for a fort, then say \emph{ghernā} again
\end{itemize}

\begin{tcolorbox}[breakable,colback=teal!4,colframe=teal!35!black,title={Before you move on}]
What does \ur{گھیرنا} mean? (\textquotedblleft{}To surround; to lay siege to\textquotedblright{}.) Say it once more.
\end{tcolorbox}
\section[\texorpdfstring{\ur{توپ} (top)}{توپ (top)}]{\ur{توپ} (top) — a cannon}

\label{lesson:UR-C309-top}



\begin{quote}
Before the new one: say the Urdu for a fort, then the Urdu for to surround.
\end{quote}

\subsection*{You'll want to know: \ur{توپ}}
\textbf{\ur{توپ}} — \emph{top} — \textquotedblleft{}a cannon\textquotedblright{}.

\begin{quote}
\ur{قلعے} \ur{پر} \ur{چار} \ur{توپیں} \ur{تھیں۔}
\end{quote}

\begin{quote}
\emph{qile par chār topeṅ thīṅ.}
\end{quote}

\begin{quote}
There were four cannons on the fort.
\end{quote}

Feminine.

\begin{cousinweb}
\textbf{\ur{توپ}} (\emph{top}) is usually traced to Turkish \emph{top}, \textquotedblleft{}a ball\textquotedblright{}, and so a cannon, which came into Urdu through Persian. It is feminine: \textbf{\ur{چار} \ur{توپیں}}.
\end{cousinweb}

\begin{grammarlens}[title={What you've built}]
One more word for telling what happened.
\end{grammarlens}

\subsection*{Your turn}
\begin{itemize}
\raggedright
  \item \emph{Say it:} \emph{top}
  \item \emph{Say it:} \emph{top}, once more
  \item \emph{Say it:} \emph{qile par chār topeṅ thīṅ.} — There were four cannons on the fort.
  \item \emph{Recall:} say the Urdu for a fort, then the Urdu for to surround, then say \emph{top} again
\end{itemize}

\begin{tcolorbox}[breakable,colback=teal!4,colframe=teal!35!black,title={Before you move on}]
What does \ur{توپ} mean? (\textquotedblleft{}A cannon\textquotedblright{}.) Say it once more.
\end{tcolorbox}
\section[\texorpdfstring{\ur{بھڑنا} (bhiṛnā)}{بھڑنا (bhiṛnā)}]{\ur{بھڑنا} (bhiṛnā) — to clash; to come to blows}

\label{lesson:UR-C309-bhirna}



\begin{quote}
Before the new one: say the Urdu for to surround, then the Urdu for a cannon.
\end{quote}

\subsection*{You'll want to know: \ur{بھڑنا}}
\textbf{\ur{بھڑنا}} — \emph{bhiṛnā} — \textquotedblleft{}to clash; to come to blows\textquotedblright{}.

\begin{quote}
\ur{دونوں} \ur{فوجیں} \ur{میدان} \ur{میں} \ur{بھڑ} \ur{گئیں۔}
\end{quote}

\begin{quote}
\emph{donoṅ faujeṅ maidān meṅ bhiṛ gaīṅ.}
\end{quote}

\begin{quote}
The two armies clashed on the field.
\end{quote}

Past (he …): \ur{وہ} \ur{بھڑ} \ur{گیا}, \emph{voh bhiṛ gayā}.

\begin{cousinweb}
\textbf{\ur{بھڑنا}} (\emph{bhiṛnā}) is inherited: CDIAL lists the Hindi cousin \emph{bhiṛnā}, \textquotedblleft{}to come close, to fight\textquotedblright{}, under a form it reconstructs (CDIAL 9490). Who you clash with takes \textbf{\ur{سے}}.
\end{cousinweb}

\begin{grammarlens}[title={What you've built}]
One more word for telling what happened.
\end{grammarlens}

\subsection*{Your turn}
\begin{itemize}
\raggedright
  \item \emph{Say it:} \emph{bhiṛnā}
  \item \emph{Say it:} \emph{bhiṛnā}, once more
  \item \emph{Say it:} \emph{donoṅ faujeṅ maidān meṅ bhiṛ gaīṅ.} — The two armies clashed on the field.
  \item \emph{Recall:} say the Urdu for to surround, then the Urdu for a cannon, then say \emph{bhiṛnā} again
\end{itemize}

\begin{tcolorbox}[breakable,colback=teal!4,colframe=teal!35!black,title={Before you move on}]
What does \ur{بھڑنا} mean? (\textquotedblleft{}To clash; to come to blows\textquotedblright{}.) Say it once more.
\end{tcolorbox}
